\documentclass[border=10pt]{standalone} 
\usepackage[utf8]{inputenc}
\usepackage{nacal}
\usepackage{pgfplots}
\usepackage{tikz-3dplot}
\usepackage{tikz} 
\usetikzlibrary{matrix,decorations.pathreplacing}
\usetikzlibrary{calc,angles,positioning,intersections,quotes,decorations.markings}
\usepackage{tkz-euclide}

\pgfplotsset{width=10cm,height=6cm}

\begin{document}
  \def\Rojo{red!50!black}
  \def\RojoOscuro{red!20!black}
  \def\RojoClaro{red!90!black}
  \def\Verde{green!50!black}
  \def\VerdeOscuro{green!30!black}
  \def\VerdeClaro{green!50!gray}
  \def\Azul{blue!50!black}
  \def\AzulOscuro{blue!35!black}
  \def\AzulClaro{blue!60!gray}

  \begin{tikzpicture}
    [scale=1.1,
    guide/.style={dashed, thick, gray},
    vector/.style={-{Stealth[length=3mm, width=2mm]},very thick},
    vector aux/.style={-{Stealth[length=2mm, width=1mm]},dashed,thick}]
    % vértices del triángulo
    \coordinate (v1) at ( 0, 0  );
    \coordinate (v2) at ( 3, 0  );
    \coordinate (v3) at (-2, 0  );
    \coordinate (v4) at (-2, 2.5);
    \coordinate (v5) at ( 0, 2.5);  
    
    % múltiplo de a
    \draw[vector aux, \AzulClaro!65!gray] (v1) -- (v3) node [below]  {$\alpha\Vect{a}$};
    % \draw[-,gray,dashed,very thick]  (v1) -- (v3);
    
    % triangulo
    \draw[vector,\Azul]  (v1) -- (v2) node [below]  {$\Vect{a}$};
    \draw[vector,\Verde] (v1) -- (v4) node [above]  {$\Vect{b}$};

    % La longiitud de C es la longitud de \Vect{b}-\Vect{a}
    \draw[guide] (v3) -- (v4);
    \draw[vector aux, \AzulClaro!65!gray,thick,dashed] (v4) -- node [below] {$-\alpha\Vect{a}$} (v5);    
    \draw[vector,\Rojo] (v1) -- (v5) node [right] {$\Vect{h}=\Vect{b}-\alpha\Vect{a}$};
    % \draw pic[<->,draw=black,angle radius=15,angle eccentricity=1.4,green!50!black,thick] {angle = v2--v1--v5};
    \tkzMarkRightAngle[size=.35,gray,very thin,fill=lightgray!70,opacity=0.5](v5,v1,v2)M
  \end{tikzpicture}
\end{document}
